\section{METODOLOGI}

Data historis harga saham yang digunakan dalam penelitian ini diperoleh secara otomatis melalui \textit{Application Programming Interface} (API) Yahoo Finance menggunakan pustaka \texttt{yfinance} pada bahasa pemrograman Python. Akuisisi data difokuskan pada saham-saham unggulan indeks IHSG dalam dua konfigurasi sistem, yaitu sistem $N=2$ yang terdiri atas saham BBCA dan TLKM, serta sistem $N=4$ yang mencakup BBCA, ADRO, SMGR, dan TLKM, pada periode Januari 2021 hingga Desember 2023. Data yang diunduh berupa harga penutupan harian yang kemudian dibersihkan dari nilai yang hilang (\textit{missing values}) dan diurutkan berdasarkan indeks waktu untuk menjaga konsistensi deret waktu. Rentang waktu tersebut dipilih untuk merepresentasikan dinamika pasar modal Indonesia pada fase pemulihan pasca-pandemi yang mencakup kondisi volatilitas yang beragam.

Implementasi seluruh alur komputasi penelitian ini dibangun di atas ekosistem pustaka Python. Pustaka \texttt{PennyLane} digunakan sebagai \textit{framework} utama untuk simulasi sirkuit kuantum, konstruksi Ansatz EfficientSU(2), serta evaluasi nilai ekspektasi Hamiltonian secara diferensiabel. Pustaka \texttt{Pandas} digunakan untuk manajemen dan praproses data deret waktu saham, sedangkan \texttt{NumPy} menangani seluruh komputasi numerik matriks kovarians, ekspektasi imbal hasil, dan kalkulasi parameter Hamiltonian Ising. Modul \texttt{Matplotlib} digunakan untuk visualisasi kurva konvergensi energi, distribusi probabilitas, dan kinerja portofolio, sementara fungsi \texttt{scipy.optimize} diterapkan sebagai pembanding dengan metode \textit{Sequential Least Squares Programming} (SLSQP) pada skenario optimasi klasik. Integrasi seluruh komponen tersebut membentuk sebuah alur kerja komputasi hibrida yang dapat dieksekusi secara penuh pada lingkungan simulator perangkat kuantum virtual tanpa meninjau derau.

Alur pengolahan data penelitian ini mencakup beberapa fase yang dapat ditinjau pada diagram alir Gambar~\ref{fig:flowchart_gt_vqe}. Fase pra-pemrosesan dimulai dari akuisisi data harga saham, transformasi ke imbal hasil sederhana dan logaritmik, kalkulasi matriks kovarians $\Sigma$ dan ekspektasi imbal hasil $\mu$, serta penentuan parameter penghindaran risiko endogen $\gamma$ menggunakan fungsi sigmoid. Fase selanjutnya adalah pencarian ekuilibrium Nash melalui kerangka EPG dan konstruksi Hamiltonian Ising berbasis QUBO. Setelah Hamiltonian terbentuk, fase optimasi kuantum dengan VQE-SPSA dijalankan secara iteratif, diikuti evaluasi entropi Von Neumann sebagai indikator konvergensi. Terakhir, fase backtesting \textit{rolling window} dilakukan untuk mengeksekusi penyeimbangan kembali portofolio secara bulanan dan mengkalkulasi metrik kinerja akhir berupa \textit{Cumulative Return}, \textit{Sharpe Ratio}, dan \textit{Maximum Drawdown}.

\begin{figure*}[t!]
    \centering
    \begin{tikzpicture}[node distance=1.3cm, every node/.style={fill=white, font=\footnotesize}, align=center]
        % Styles
        \tikzstyle{startstop} = [ellipse, minimum width=2.5cm, minimum height=0.6cm, text centered, draw=black, fill=white]
        \tikzstyle{process} = [rectangle, minimum width=3.2cm, minimum height=0.8cm, text centered, draw=black, fill=white]
        \tikzstyle{decision} = [diamond, minimum width=2.5cm, minimum height=1cm, text centered, draw=black, fill=white, aspect=2]
        \tikzstyle{connector} = [circle, draw=black, minimum size=0.5cm, text centered, inner sep=2pt]
        \tikzstyle{arrow} = [thick,->,>=stealth]
        
        % KOLOM 1: Pra-pemrosesan & Parameterisasi
        \node (start) [startstop] {Mulai};
        \node (data) [process, below of=start, yshift=-0.5cm] {Akuisisi Data \& \\ Setup Jendela Waktu};
        \node (param) [process, below of=data, yshift=-0.5cm] {Hitung $\mu$, $\Sigma$ \& \\ penghindaran risiko ($\gamma$)};
        \node (Nash) [process, below of=param, yshift=-0.5cm] {Pencarian ekuilibrium Nash \\ dengan Kerangka EPG};
        \node (hamil) [process, below of=Nash, yshift=-0.5cm] {Formulasi QUBO \& \\ Hamiltonian ($h_i, J_{ij}$)};
        \node (c1a) [connector, below of=hamil, yshift=-0.5cm] {1};
        \node (c3b) [connector, left of=data, xshift=-1cm] {3};
        
        % KOLOM 2: Optimasi Kuantum (VQE)
        \node (c1b) [connector, right of=start, xshift=4.0cm] {1};
        \node (vqe_init) [process, below of=c1b, yshift=-0.2cm] {Inisialisasi State Awal \\ $|0\dots0\rangle \to |\text{bitstring}_{NE}\rangle$ \\ \& Parameter $\theta$};
        \node (spsa) [process, below of=vqe_init, yshift=-0.5cm] {Optimasi SPSA dan \\ Pengukuran Energi};
        \node (entropy) [process, below of=spsa, yshift=-0.5cm] {Kalkulasi Entropi dari \\ Konfigurasi Probabilitas};
        \node (conv) [decision, below of=entropy, yshift=-0.5cm] {Konvergen?};
        \node (layer_eval) [process, below of=conv, yshift=-0.6cm] {Simpan Energi $E_d$ \\ pada kedalaman $d$};
        \node (layer_dec) [decision, below of=layer_eval, yshift=-0.5cm] {$d < D_{max}$?};
        \node (c2a) [connector, below of=layer_dec, yshift=-0.7cm] {2};
        
        % KOLOM 3: Evaluasi & Backtesting
        \node (c2b) [connector, right of=c1b, xshift=4.0cm] {2};
        \node (best_state) [process, below of=c2b, yshift=-0.2cm] {Pilih Bitstring \\ dengan $E_{min}$};
        \node (rebal) [process, below of=best_state, yshift=-0.5cm] {Eksekusi penyeimbangan kembali \\ Portofolio};
        \node (end_dec) [decision, below of=rebal, yshift=-1cm] {Akhir \\ Periode?};
        \node (eval) [process, below of=end_dec, yshift=-1cm] {Kalkulasi Metrik \\ Kinerja Akhir};
        \node (stop) [startstop, below of=eval, yshift=-0.5cm] {Selesai};
        \node (c3a) [connector, right of=end_dec, xshift=2.1cm] {3};

        % Alur Kolom 1
        \draw [arrow] (start) -- (data);
        \draw [arrow] (data) -- (param);
        \draw [arrow] (param) -- (Nash);
        \draw [arrow] (Nash) -- (hamil);
        \draw [arrow] (hamil) -- (c1a);
        \draw [arrow] (c3b) -- (data);

        % Alur Kolom 2
        \draw [arrow] (c1b) -- (vqe_init);
        \draw [arrow] (vqe_init) -- (spsa);
        \draw [arrow] (spsa) -- (entropy);
        \draw [arrow] (entropy) -- (conv);
        \draw [arrow] (conv) -- node[anchor=east] {Ya} (layer_eval);
        \draw [arrow] (conv.west) -- ++(-0.6,0) |- (spsa.west) node[pos=0.1, anchor=west] {Tidak};
        \draw [arrow] (layer_eval) -- (layer_dec);
        \draw [arrow] (layer_dec) -- node[anchor=east] {Tidak} (c2a);
        \draw [arrow] (layer_dec.east) -- ++(0.7,0) |- (vqe_init.east) node[pos=0, anchor=west] {Ya (Tambah \\ kedalaman)};


        % Alur Kolom 3
        \draw [arrow] (c2b) -- (best_state);
        \draw [arrow] (best_state) -- (rebal);
        \draw [arrow] (rebal) -- (end_dec);
        \draw [arrow] (end_dec) -- node[anchor=east] {Ya} (eval);
        \draw [arrow] (eval) -- (stop);
        \draw [arrow] (end_dec.east) -- (c3a.west) node[midway, anchor=south] {Tidak};

    \end{tikzpicture}
    \caption{Diagram Alir Metodologi Optimasi Portofolio Berbasis GT-VQE.}
    \label{fig:flowchart_gt_vqe}
\end{figure*}

\subsection{Fungsi Objektif Markowitz dan Formulasi Biner}

Secara matematis, awal dari alur metodologi penelitian ini dimulai dari fungsi objektif portofolio \textit{Mean-Variance} Markowitz. Model Markowitz berupaya meminimalkan varians portofolio sambil mempertahankan ekspektasi imbal hasil pada level tertentu, yang dalam formulasi Lagrangian dapat dinyatakan sebagai
\begin{equation}
    f(\mathbf{w}) = \frac{\lambda}{2} \sum_{i,j} \sigma_{ij} \omega_i \omega_j - \sum_i \mu_i \omega_i,
\end{equation}
dengan $\omega_i$ adalah vektor bobot aset kontinu, $\mu_i$ adalah ekspektasi imbal hasil aset ke-$i$, $\sigma_{ij}$ adalah elemen matriks kovarians, dan $\lambda$ adalah parameter penghindaran risiko. Karena penyelesaian langsung pada ruang bobot kontinu menjadi \textit{NP-Hard} ketika batasan kardinalitas diberlakukan, bobot kontinu $\omega_i$ didiskritisasi menjadi variabel biner $x_i \in \{0, 1\}$ dengan bobot portofolio $\omega_i = x_i/K$, sehingga fungsi objektif Lagrangian berubah menjadi bentuk biner
\begin{equation}
    f(\mathbf{x}) = \frac{\lambda}{2K^2} \sum_{i,j} \sigma_{ij} x_i x_j - \frac{1}{K}\sum_i \mu_i x_i,
    \label{eq:lagrangian_markowitz_binary}
\end{equation}
di mana $K$ adalah jumlah aset yang dipilih ke dalam portofolio. Representasi biner inilah yang menjadi titik masuk transformasi ke dalam kerangka \textit{Exact Potential Game} dan selanjutnya ke dalam formulasi QUBO untuk komputasi kuantum.

\subsection{Fungsi Potensial Global dan Keseimbangan Nash}

Tahapan pertama dalam menyusun struktur komputasi alokasi aset adalah mentransformasikan persamaan Lagrangian Markowitz biner standar menjadi sebuah model \textit{Exact Potential Game} (EPG) \cite{monderer_potential_1996}. Transformasi matematis ini dilakukan dengan menyubstitusikan variabel bobot kontinu ke dalam representasi diskret $x_i \in \{0, 1\}$, yang mengindikasikan status keterpilihan suatu instrumen aset di dalam konfigurasi portofolio. Integrasi faktor penghindaran risiko, yang disimbolkan dengan $\gamma$, diinjeksikan ke dalam model permainan untuk mengendalikan sensitivitas matriks kovarians, yang merepresentasikan ukuran interaksi antara aset di dalam sistem pasar yang kompetitif. Dinamika utilitas yang beroperasi pada setiap aset tunggal akibat interaksi strategis tersebut kemudian dirangkum sepenuhnya menjadi sebuah fungsi skalar tunggal yang dikenal sebagai fungsi potensial global ($\Phi(\mathbf{x})$). Fungsi potensial ini secara eksplisit merangkum seluruh kalkulasi imbal hasil historis ($\mu$) serta menyeimbangkannya dengan tingkat risiko volatilitas absolut dari keseluruhan kombinasi strategis agen yang saling bersinggungan. Secara keseluruhan, fungsi ini memiliki bentuk
\begin{equation}
\label{eq:global_potential_function}
\Phi(\mathbf{x}) = \sum_{l=1}^N \mu_l \frac{x_l}{K} - \frac{\gamma}{2}\sum_{i=1}^N\sum_{j=1}^N \sigma_{ij} \frac{x_i}{K} \frac{x_j}{K}.
\end{equation}

Fungsi potensial global tersebut menjamin bahwa model alokasi investasi yang diusulkan memiliki lintasan konvergensi matematis yang stabil dan deterministik \cite{jain_multi-portfolio_2012}. Setiap transisi strategi dari agen individu yang berhasil meningkatkan utilitas lokalnya akan serta merta berkontribusi pada peningkatan nilai fungsi potensial makroskopis sistem dengan kuantitas deviasi yang presisi ekuivalen ($\Delta u_i = \Delta \Phi$). Evolusi pergerakan dinamis ini akan otomatis terhenti ketika sistem mencapai konfigurasi optimum lokal, yang direpresentasikan oleh ketiadaan perubahan strategis alternatif yang mampu memberikan peningkatan nilai fungsi lebih lanjut. Kondisi stabil puncak dari evolusi fungsi potensial ini dalam literatur fisika sosial direpresentasikan sebagai ekuilibrium \textit{Nash}, yang posisinya akan dideteksi secara heuristik menggunakan algoritma \textit{Sequential Best Response} sebelum diumpankan ke fase evaluasi kuantum. Keseimbangan strategis yang berhasil diidentifikasi akan direkonstruksi kembali ke dalam bentuk vektor biner, berfungsi sebagai basis parameter pemicu awal yang esensial dalam kerangka operasional algoritma VQE secara keseluruhan.

\subsection{Transformasi Hamiltonian Ising}

Untuk memungkinkan pemrosesan algoritma VQE pada infrastruktur perangkat keras kuantum, bentuk fungsi potensial EPG harus diproyeksikan ke dalam struktur matriks Hamiltonian \textit{Ising} \cite{lucas_ising_2014}. Tahapan ini diawali dengan mengonversi rumusan maksimisasi potensial global menjadi fungsi minimisasi energi terhadap domain keputusan ($E(\mathbf{x}) = -\Phi(\mathbf{x})$). Hasil persamaan energi tersebut kemudian dikembangkan lebih lanjut ke dalam model \textit{Quadratic Unconstrained Binary Optimization} (QUBO) dengan menambahkan fungsi penalti kuadratik $P(\mathbf{x}) = A\left(\sum_{i=1}^N x_i - K\right)^2$. Variabel penalti tambahan $A$ digunakan untuk memastikan bahwa matriks solusi algoritma memenuhi batasan seleksi kardinalitas berjumlah $K$ aset sesuai dengan yang diinginkan oleh pengguna sebagai investor dan dipenuhi di awal eksperimen. Sintesis fungsional dari variabel energi potensial dan fungsi penalti tersebut memformulasikan koefisien matriks relasional yang komprehensif pada variabel $Q_{ii}$ serta elemen korelatif interaksi $Q_{ij}$.

Setelah kedua elemen $Q$ didapatkan, dilakukan transisi final ke ranah komputasi kuantum dilaksanakan melalui pemetaan variabel keputusan relasional dari dominan representasi biner menuju domain mekanika \textit{spin} yang didefinisikan dengan vektor spektral $s_i \in \{-1, 1\}$. Transisi tersebut menggunakan pemetaan affine, yang menyatakan hubungan $x_i = (1 - s_i)/2$. pemetaan tersebut disubstitusi pada setiap koefisien QUBO yang ada untuk mengekstraksi nilai parameter operator uniter lokal Pauli-Z \cite{nakahara_quantum_2008}. Ekstraksi tersebut akan menghasilkan persamaan analitik untuk komponen medan bias lokal ($h_i$) pada setiap \textit{qubit} individual dan matriks kopling keterbelitan ($J_{ij}$) antar sepasang interaksi \textit{qubit} spesifik di dalam jaringan sebagai Hamiltonian Ising berbentuk
\begin{equation}
\label{eq:ising}
\hat{\mathcal{H}} = - \sum_{i=1}^N h_i \hat{Z}_i - \sum_{i<j}^N J_{ij} \hat{Z}_i \hat{Z}_j + C
\end{equation}
Komponen bias linier $h_i$ menganalogikan representasi sentimen fundamental pasar, sedangkan distribusi nilai kopling $J_{ij}$ memetakan fungsi korelasi harga dari matriks risiko. Hamiltonian Ising akhir tersebut selanjutnya digunakan sebagai observabel fungsi objektif dari VQE untuk mencari konfigurasi probabilitas status minimum dasar (\textit{ground state}) sistem yang paling mendekati solusi optimal menggunakan Ansatz dan optimasi yang telah ditentukan.

\subsection{Backtesting dan Benchmarking}

Hasil konfigurasi pemilihan aset dari sistem optimasi hibrida yang telah dieksekusi, dapat dievaluasi melalui mekanisme \textit{backtesting} menggunakan konsep \textit{rolling window}. Data deret waktu pergerakan harga saham-saham dari indeks IHSG (yakni kombinasi BBCA, ADRO, TLKM, SMGR) selama 6 bulan diekstrak secara berurutan untuk mendefinisikan matriks imbal hasil ($\mu$) dan matriks kovarians ($\Sigma$) pada tiap periode 1 bulan. Setiap periode, dilakukan kalkulasi ulang fungsi matriks \textit{Ising} baru sambil tetap menggunakan metode injeksi ekuilibrium Nash sebagai \textit{warm-start} dari mekanisme \textit{Exact Potential Game}. Sinyal investasi yang didapatkan berupa rekomendasi biner yang terdiri dari 0 yang berarti tidak beli dan 1 yang berarti beli. Konfigurasi tersebut didapatkan dari kejatuhan \textit{ground state} akhir algoritma VQE yang dieksekusi secara bertahap untuk mengatur perimbangan komposisi persentase dana dari nilai modal aktual (\textit{rebalancing}). Hasil kinerja performa akhir kemudian dihimpun ke dalam profil \textit{return} sehingga kemampuan profitabilitas portofolio kuantum ini tervalidasi.

Validitas keandalan optimasi portofolio GT-VQE juga membutuhkan standarisasi performa (\textit{benchmarking}) eksternal melalui analisis komparatif yang ketat. Proses kalibrasi referensi menggunakan pengaplikasian algoritma klasik mapan berjenis \textit{Sequential Least Squares Programming} (SLSQP), yang merupakan pengoptimal fungsi kuadratik standar pada model Markowitz linear. Evaluasi SLSQP dilakukan secara berurutan pada siklus \textit{backtesting} dari data historis untuk menyajikan dasar referensi yang transparan. Selain model konvensional optimasi tersebut, komparasi yang digunakan juga meninjau kinerja pergerakan aset konstan jika investor membeli semua aset secara merata dari awal periode hingga akhir periode (\textit{Equal Weight Portfolio}) dengan vektor matriks pembagian fraksi sebesar $\omega_i = 1/N$ \cite{demiguel_optimal_2009}. Metrik lain yang digunakan yaitu \textit{Cumulative Return}, \textit{Sharpe Ratio}, dan \textit{Maximum Drawdown} yang untuk membuktikan bahwa dimensi tinggi menggunakan pendekatan kuantum hibrida akan lebih andal daripada menggunakan algoritma klasik terhadap pasar saat ini.
